Without code context the model never produced an applicable edit: none of its 19 parsed edits matched the file,
because it quoted code it had not been shown.

With the BM25 whole-file context it produced \EtoeBmPatches{} patches and resolved none. With the goblin packet it
produced \EtoeGoblinPatches{} patches and resolved one task (pytest-7373, gold-valid). The packet condition used
\EtoeGoblinInput{} provider input tokens per task against \EtoeBmInput{}, a 3.3-fold reduction, at about the same
output length (336 against 370 tokens) and generation time (49 against 48 s).

The sample is small and the 4B model is weak, so one resolved task against none is not a significant difference.
What the run does establish:
\begin{itemize}
\item the full agentic path works end to end: Goose GDK app, local MCP router, packet, model, patch, official
harness;
\item the compact packet gives a model enough exact code to write applicable patches at under a third of the
standard retrieval context;
\item resolution, the headline SWE-bench metric, is now measured for this pipeline and can be re-run with a stronger
model on the same harness.
\end{itemize}
